\section{Empirical design and inference}
\subsection{The estimand and its causal boundary}

Let $r_{it}$ denote security $i$'s daily return and $r_{mt}$ the OMXSBGI return, both expressed in percentage points. The market-adjusted cumulative return over a trading-day window $[a,b]$ is
\begin{equation}
 CAR_i[a,b]=\sum_{t=a}^{b}(r_{it}-r_{mt}).
 \label{eq:car}
\end{equation}
The primary window is $[AD_i-1,ED_i+1]$. Short announcement and implementation windows are $[AD_i-1,AD_i+1]$ and $[ED_i-1,ED_i+1]$. These quantities are arithmetic sums of daily excess returns, not compounded buy-and-hold abnormal returns. The short-window means need not add to the primary-window mean because the primary window also includes intervening days, and because completeness requirements can select different observations.

For event type $s$ and period $p$, the principal parameter is the mean of this measured CAR among usable events of that type and period. It is not the average treatment effect of index membership. To identify such an effect, the benchmark-adjusted counterfactual return would have to remove relevant common movements, selection into membership would need to be conditionally unrelated to the remaining counterfactual outcome, and coincident information would need to be excluded or accounted for. The present design does not establish those conditions.

An illustrative decomposition clarifies the problem. Suppose an observed event return can be written as
\begin{equation}
 CAR_i=\lambda_i Q_i+I_i+B_i+U_i,
 \label{eq:decomp}
\end{equation}
where $Q_i$ is unexpected benchmark-related demand, $\lambda_i$ its price response, $I_i$ is contemporaneous information, $B_i$ is selection-related expected return, and $U_i$ collects remaining benchmark and measurement error. Equation~\eqref{eq:decomp} is an interpretive accounting device, not an estimated structural model. A change in mean CAR can reflect any component, including covariance between demand and its price response. Neither index membership nor a positive holdings change separately identifies $\lambda_i$.

This framework yields discriminating questions rather than mutually exclusive mechanism tests. A stronger pre-announcement return is consistent with anticipation but also with selection on earlier performance. Lower quoted spreads are consistent with cheaper trading but do not identify the supply curve faced by benchmark traders. Positive post-event returns can reflect incomplete reversal, continued information or a misspecified benchmark. We evaluate these implications while retaining the alternatives.

\subsection{Period contrasts and observable composition}

The unadjusted specification regresses CAR on four mutually exclusive period indicators. Each coefficient is the event-weighted period mean. Linear contrasts test the difference between the earliest and latest blocks, the last two blocks, and the two decades. A separate regression on calendar year summarizes a linear trend without imposing the four bins. The trend is a descriptive projection over an uneven historical sample, not evidence of a constant structural decline.

The composition specification adds centered measures of abnormal turnover, relative market capitalization, residual return variance and analyst coverage. It is estimated separately for additions and deletions. The inherited variable construction uses information near AD, and some variables include AD-1, which also lies in the outcome window. This timing overlap is another reason to interpret the exercise as an adjustment rather than a clean conditioning strategy for causal identification. We report an unadjusted regression on the identical complete-case sample beside the adjusted regression. A change between full-sample and adjusted estimates otherwise mixes sample selection with covariate adjustment.

The comparison-security regression adds review fixed effects, log capitalization and a pre-announcement return measure to an event indicator. Nonmember comparison securities for additions pass a size screen based on the smallest addition; member comparisons for deletions pass the corresponding screen based on the largest deletion. These screens are not the Nasdaq assignment rule. They neither demonstrate local random assignment nor guarantee comparable latent fundamentals. The comparison coefficient is therefore a conditional return difference. Moreover, the inherited pre-return control includes AD-1 and overlaps the outcome by one day. The regression remains a secondary descriptive diagnostic until that overlap and historical eligibility construction are independently rebuilt.

\subsection{Demand scaled returns}

Let $s_i=1$ for additions and $s_i=-1$ for deletions. We sign-adjust both returns and holdings changes so that positive values correspond to index-directed movements. For a covered fund group $g$, define
\begin{equation}
 D_{ig}=100s_i\frac{\sum_{f\in g}\Delta H_{if}}{S_i},\qquad
 Y_i=s_iCAR_i[AD_i-1,ED_i+1],\qquad
 \widehat M_{pg}=\frac{\overline{Y}_{pg}}{\overline{D}_{pg}},
 \label{eq:ratio}
\end{equation}
where $\Delta H_{if}$ is the change in covered fund $f$'s reported shares and $S_i$ is shares outstanding. Both averages use the same non-split complete cases. The ratio measures percentage points of sign-adjusted CAR per percentage point of measured ownership change. It is a ratio of means, not a mean of event-level ratios.

The distinction matters statistically and economically. A small mean denominator can make the ratio large even when its numerator is imprecise. If the covered holdings measure is $D^{obs}=\kappa Q+\eta$, where $\kappa$ captures coverage and $\eta$ measurement error or unrelated trading, the ratio is not invariant to changes in coverage. Even with negligible $\eta$, a changing $\kappa$ changes the observed ratio without a change in the underlying price response. Fund assets, disclosed holdings and actual implementation orders are therefore not interchangeable denominator measures.

The inherited analysis reports percentile intervals from resampling review clusters. We retain those outputs in the audit package but emphasize a Fieller sensitivity calculation that uses the joint covariance of the two means. A candidate ratio $m$ is retained when
\begin{equation}
 (\overline Y-m\overline D)^2\leq
 t_{G-1,0.975}^{2}
 \left[\widehat V_{YY}-2m\widehat V_{YD}+m^2\widehat V_{DD}\right].
 \label{eq:fieller}
\end{equation}
Here $\widehat V$ is the review-clustered CR1 covariance matrix of $(\overline Y,\overline D)$. The set can be bounded, consist of two unbounded rays, or include every real value. Such outcomes are information about weak denominator precision, not computational failures to be discarded. With very few reviews, the $t$ approximation is itself uncertain; the Fieller analysis is an informative sensitivity calculation rather than a claim of exact coverage \citep{fieller1954}.

\subsection{Clustering and finite sample uncertainty}

The main return tables use the inherited restricted wild cluster bootstrap. For each null hypothesis, the code imposes the coefficient restriction, multiplies residuals by a common random weight within each review, re-estimates the model and compares studentized statistics. It uses 9,999 draws, Rademacher weights when at least 12 clusters are available and Webb weights otherwise. Confidence intervals are obtained by numerical test inversion. The review is the natural first clustering unit because several securities share an announcement and implementation cycle.

This procedure does not make reviews independent by construction. A firm can reappear in multiple reviews, macroeconomic conditions can persist, and a period comparison may have only a few influential reviews behind one side of the contrast. We therefore supplement the main results with review jackknife uncertainty and review-by-security two-way clustering, following the dependence concerns discussed by \citet{petersen2009,cameron2011,mackinnon2023}. Equal-review means answer a different weighting question and are reported as such. Leave-one-review-out ranges show influence; they are not confidence intervals.

The rerun audit also reveals a numerical limitation in some auxiliary bootstrap inversions. The original routine searches outward from the point estimate, stops at a finite search boundary, and returns a single interval. In sparse designs, that need not characterize a potentially irregular acceptance set. Floating-point differences can change some of those endpoints substantially. We do not rely on unstable early-period auxiliary ratio intervals. Composition-adjustment intervals in the manuscript use analytical CR1 uncertainty, explicitly labeled, and the ratio table reports Equation~\eqref{eq:fieller}. Main return-bootstrap intervals remain the frozen pointwise results, whose rerun differences are small and documented.

\subsection{Multiplicity, power and economic disappearance}

The research record contains many periods, windows, screens and earlier exploratory analyses. It would be inaccurate to describe all reported tests as fixed before the data were observed. We treat the inherited analysis as exploratory and avoid selecting the preferred specification by its significance. As a limited diagnostic, Appendix~\ref{app:multi} applies Holm adjustment to a declared family of ten inherited return tests. This does not account for all historical researcher choices and cannot retroactively create a confirmatory design \citep{holm1979}.

The archived minimum-detectable-effect calculation is
\begin{equation}
 MDE\approx\left(t_{G-1,0.975}+t_{G-1,0.80}\right)\widehat{SE}.
 \label{eq:mde}
\end{equation}
This is a plug-in analytical approximation for a two-sided 5\% test at 80\% power. It is not the exact power of the implemented wild bootstrap, nor does it account for the full specification search. We use it to describe the scale of differences that the sample can resolve, rather than calculating ``observed power'' as a new argument for the estimated effect.

Disappearance requires more than failure to reject zero. For a transparent audit benchmark, we ask whether a 90\% analytical confidence interval for the latest-period mean lies wholly within $[-1,1]$ percentage points, corresponding to two one-sided 5\% equivalence tests. The margin is chosen for reporting after the old results were known, not as a preregistered or universally valid economic threshold. Failure of this check means that this dataset does not establish practical equivalence at that margin; it does not prove that the true effect exceeds it.

\subsection{A fixed cohort horizon extension}

The new extension divides the cumulative return into non-overlapping intervals:
\begin{equation}
 CAR_i[AD_i-1,ED_i+20]
 =CAR_i[AD_i-1,ED_i-1]+CAR_i[ED_i,ED_i+20].
 \label{eq:horizon}
\end{equation}
All three quantities are calculated on the same complete-case cohort. This avoids combining a short-window mean with a post-event mean based on different securities, or counting ED and ED+1 twice. We report all four five-year blocks and both event types in the machine-readable audit, with the later decade and latest block emphasized in the main table.

The extension uses independently coded analytical CR1 intervals and a separate restricted wild-bootstrap test of the mean with 49,999 draws. Simulated $p$-values use $(k+1)/(B+1)$, where $k$ is the number of exceedances. The bootstrap $p$-value and analytical interval are deliberately labeled as different summaries. The extension was specified after inherited outcomes were known; additional exclusions were examined after viewing the new results. A nonzero cumulative mean rejects a particular zero-mean finite-horizon hypothesis under the maintained return model. It does not identify an index-induced permanent level shift.
